\documentclass{article}
\title{Trust in Local Government After a Municipal Merger}
\date{}
\begin{document}
\maketitle

\begin{abstract}
We examined how residents' trust in local government changed after two municipalities merged. In interviews with 24 residents, trust fell where the merger moved services away. A survey of 412 residents showed lower trust among those who lost a local office.
\end{abstract}

\noindent Keywords: political trust, municipal merger, local government, sociology

\section{Introduction}

A municipal merger moves public services away from some residents. Residents who lose a local office may feel less political trust, the belief that a government acts in the interest of its citizens (Abara, 2017).

The study followed residents of two towns that merged in 2019.

Merged municipalities now serve 18 percent of residents in the country (Statistics Office, 2022).

Interviews explored how residents understood the merger, and a survey measured trust.

\section{Interview Method}

Semi-structured interviews were held with 24 residents (N = 24). Participants were recruited through the public library.

Residents were eligible if they had lived in either town since before the merger.

The interview guide asked about services and about contact with the council.

The researcher conducted the interviews in the participants' homes.

Each interview lasted 45 minutes. Interviews were recorded and transcribed.

Each transcript was coded twice for themes of distance and voice.

\section{Interview Results}

Fifteen of the 24 participants said that the new council felt distant. They described the loss as a loss of voice: the council no longer listened to them.

The merger brought better services in the view of six participants. Each of them lived within 5 km of the new town hall, which may explain the view.

Without a local office, participants faced a trip of 20 to 40 minutes to the new town hall. Four of them stopped visiting, which suggests that distance cut their contact with the council.

The three participants who kept an office nearby reported no change in their contact with the council. A nearby office may therefore have preserved that contact.

One participant said that the council "stopped listening." The quotation illustrates the loss of voice that 15 participants described.

\section{Survey Method}

A mailed survey tested whether losing a local office lowered trust. We operationalized trust as agreement with five statements on a 5-point scale.

Place attachment was measured in the survey.

The survey was mailed to 1,200 residents, and 412 replied. A response rate of 34 percent left room for non-response bias.

People who moved after the merger were missing from the sample, because only residents with a registered address could be reached.

Losing an office was coded from the residents' registered address.

The survey had 18 items and took 10 minutes to complete. It was tested on 20 residents before mailing.

\section{Survey Results}

Respondents who had lost a local office reported lower trust than the others (mean 2.9 against 3.6, p below .001). The difference was 0.7 points on the 5-point scale. Those who had lost an office thus fell just below the scale's neutral point, and the others above it.

The respondents were a mixed bag in age and occupation, with ages from 22 to 87 years and 31 occupations.

Trust rose with years of residence (r = .18, p below .001) and showed no link to age (r = .04, p = .41). Years of residence thus related more closely to trust than age did.

Trust was lowest among residents who lived more than 10 km from the new town hall (mean 2.6). The survey result agreed with the participants' accounts of contact lost to distance.

Response rates were similar in the two towns, at 33 and 35 percent, so neither town dominated the sample.

\section{Discussion}

The interviews and the survey suggest that losing a local office lowers trust. The study adds evidence from two merged towns to the claim of Abara (2017).

Participants were recruited at the public library, so they may differ from residents who never visit it.

The pattern may also reflect residents who already distrusted the merger. A survey before the merger would separate the two explanations.

The interviews suggest a mechanism; the survey shows its size.

Local offices give residents a place to be heard (Abara, 2017). A longer trip to the town hall may reduce contact with the council.

The loss appears to concern voice, not service quality, since six participants described better services.

We suggest that councils consider keeping a local office open after a merger.

\section{Conclusion}

Losing a local office may lower residents' trust in a merged municipality. Future studies should follow residents before and after a merger.


\end{document}
